\documentclass[a4paper,12pt]{article}
\usepackage[spanish,activeacute]{babel}
\usepackage[latin1]{inputenc}

%% Para las columnas
\usepackage{multicol}
\usepackage{amssymb,amsmath,eepic,fancyhdr}
\usepackage{latexsym}
\usepackage{datenumber}


%% Para numeraci?n autom?tica
\newcounter{nroproblema}
\setcounter{nroproblema}{1}
\newcounter{nroapartado}
\setcounter{nroapartado}{1}

\newcommand{\n}{%
    \vspace{.1in}\noindent{\bf \arabic{nroproblema}.\ }%
    \addtocounter{nroproblema}{1}%
    \setcounter{nroapartado}{1}%
    }%

\newcommand{\nn}{%
    \vspace{.1in}\noindent \hskip.4cm {\bf \arabic{nroproblema}.\ }%
    \addtocounter{nroproblema}{1}%
    \setcounter{nroapartado}{1}%
    }%


\newcommand{\ap}{\vspace*{0.2cm}\hspace*{0.2cm}%
   {\bf \alph{nroapartado})}
   \addtocounter{nroapartado}{1}%
   }

 
%marcos by Marco
\newcommand{\pd}[1]{\frac{\partial}{\partial #1}} %\pd{x}
\newcommand{\NN}{\ensuremath{\mathbb N}}
\newcommand{\ZZ}{\ensuremath{\mathbb Z}}
\newcommand{\CC}{\ensuremath{\mathbb C}}
\newcommand{\RR}{\ensuremath{\mathbb R}}
\newcommand{\TT}{\ensuremath{\mathbb T}}
\newcommand{\g}{\ensuremath{\frak{g}}}
\newcommand{\h}{\ensuremath{\frak{h}}}
\newcommand{\ft}{\ensuremath{\frak{t}}}
\newcommand{\cB}{\mathcal{B}}
\newcommand{\cN}{\mathcal{N}}
\newcommand{\cL}{\mathcal{L}}
\newcommand{\cD}{\mathcal{D}}

 

%% m�rgenes

% Anterior:

%\setlength{\unitlength}{1mm} \addtolength{\voffset}{-15mm}
%\addtolength{\textheight}{30mm} \addtolength{\hoffset}{-20mm}
%\addtolength{\textwidth}{45mm}

\setlength{\unitlength}{1mm} \addtolength{\voffset}{-22mm}
\addtolength{\textheight}{40mm} \addtolength{\hoffset}{-22mm}
\addtolength{\textwidth}{46mm}




\begin{document}

\noindent
{\sffamily\large Marco Zambon 
\hfill  Master UAM, 2013-14  \\Geometr\'ia simpl\'ectica y de Poisson}


 

\begin{center}\bfseries\large\textsf{Hoja 1    \vspace{-2mm}}
\end{center}


%\begin{flushright}
%1${}^o$ de grado en F�sica, curso 2013/14
%\today
%\end{flushright}
%\begin{center}{\bf {\Large \'ALGEBRA I. HOJA 4}}\end{center}
%\vskip.35cm
 
 
 
 
 
 
 
 
 
\n 
Sea $(V,\Omega)$ un  espacio vectorial simpl\'ectico de dimensi\'on $2n$. Sean $W$ y $Z$ subespacios vectoriales de $V$. Demuestra:


\ap ($W^{\Omega})^{\Omega}=W$

\ap $W$ isotropico $\Rightarrow$ $dim(W)\le n$

\ap $W$ coisotropico $\Rightarrow$ $dim(W)\ge n$ 

\ap $W$ lagrangiano $\Rightarrow$ $dim(W)= n$

\ap si  $dim(W)=2n-1$, entonces $W$ es coisotropico

\ap $W^{\Omega}\cap Z^{\Omega}=(W+Z)^{\Omega}$ 


\ap $(W\cap Z)^{\Omega}=W^{\Omega}+Z^{\Omega}$.

\n Sea $(V,\Omega)$ un  espacio vectorial simpl\'ectico de dimensi\'on $2n$. 

\ap Si $L$ es un subespacio lagrangiano de $V$, entonces cualquier subespacio $Z$ de $V$ que contiene $L$ es coisotropico.

\ap Sea $W$  un subespacio coisotropico de $V$. Entonces existe un subespacio lagrangiano $L$ de $(V,\Omega)$ con $L\subset W$.

\noindent{\bf Consejo para b):} Utiliza el espacio vectorial simpl\'ectico $W/W^{\Omega}$.
 

\n 
Sean $(V,\Omega_V)$ y $(W,\Omega_W)$ espacios vectoriales simpl\'ecticos.
Sea $\Phi \colon V \to W$ una aplicaci\'on  lineal tal que $\Phi^*\Omega_W=\Omega_V$.

\ap Demuestra que $dim(V)\le dim(W)$

\ap Demuestra que $grafo(\Phi):=\{(v,\Phi(v)):v\in V\}$ es un subespacio isotropico de $(V\times W,\omega)$,
donde la forma simpl\'ectica lineal $\omega$ est\'a definida por $$\omega((v_1,w_1),(v_2,w_2))=\Omega_V(v_1,v_2)-\Omega_W(w_1,w_2).$$

\noindent \textbf{Recuerda:} $\Phi^*\Omega_W$ est\'a definida por  $(\Phi^*\Omega_W)(v_1,v_2)=\Omega_W(\Phi(v_1),\Phi(v_2))$ para todo $v_1,v_2\in V$.

\n Sea $V$ un espacio vectorial real, y $J\colon V\to V$ una estructura compleja. (En otras palabras,  $(V,J)$ es un espacio vectorial complejo).

\ap si $\Omega$ es una forma simpl\'ectica en $V$ tal que $J$ y $\Omega$ son compatibles, y definiendo $g:=\Omega(\cdot,J\cdot)$, entonces $h:=g-i\Omega$ es un producto hermitico en $(V,J)$

\ap reciprocamente: si $h$ es un producto hermitico en $(V,J)$,
entonces $\Omega(v,w):=-Im(h(v,w))$ define una forma simpl\'ectica $\Omega$ en $V$ compatible con $J$, y
$g(v,w):=Re(h(v,w))$ es
el producto escalar asociado (es decir, $g=\Omega(\cdot,J\cdot)$).
Aqu\'i Re y Im denotan la parte real y imaginaria.

\noindent{\bf Recuerda:} Un producto hermitico en $(V,J)$ es una aplicaci\'on
$h\colon V\times V\to \CC$ tal que:
\begin{itemize}
\item $h$ es $\CC$-lineal en la primera entrada y $\CC$-antilineal en la segunda
\item $h(v,w)=\overline{h(w,v)}$ para todo $v,w\in V$
\item $h(v,v)>0$ para todo $v\neq 0$.
\end{itemize}
 
 
 \n Sea $(V,\Omega)$ un  espacio vectorial simpl\'ectico de dimensi\'on $2n$, $J$ una estructura compleja compatible, y $g$ el producto escalar  $g:=\Omega(\cdot,J\cdot)$. Demuestra: existe un isomorfismo $V\cong \RR^{2n}$ que identifica $\Omega, J, g$ con la forma simpl\'ectica can\'onica, la estructura compleja can\'onica y el producto escalar can\'onico en $\RR^{2n}=\CC^n$. 
 
 
 \n Ejercicio 2 en p\'agina 72 del libro de Cannas da Silva (sobre $Sp(2n), O(2n), GL(n,\CC), U(n)$). 
 
\noindent{\bf Recuerda:} $U(n)$ consiste de las matrices que preservan el producto hermitico standard en $\CC^n$. 
\end{document}
